\chapter*{\texorpdfstring{{\tnr\bfseries\fontsize{14pt}{20pt}\selectfont ABSTRACT}}{ABSTRACT}}
\thispagestyle{thesis-front}
\markboth{成都信息工程大学学士学位论文}{成都信息工程大学学士学位论文}

\begin{thesisenabstract}
Real magnetic resonance imaging (MRI) data is difficult to use directly in teaching demonstrations and generative model experiments because of acquisition cost, privacy protection, and ethical review. A single two-dimensional slice also gives only limited insight into the spatial structure of the brain. Based on this practical need, this thesis designs and implements an MRI medical image pseudo-generation platform built on 3D Gaussian Splatting (3DGS) and StyleGAN2. The IXI brain MRI dataset is used as the data source. The preprocessing stage includes Neuroimaging Informatics Technology Initiative (NIfTI) volume slicing, intensity normalization, resolution standardization, low-information slice filtering, and StyleGAN2-ADA dataset construction. A two-dimensional MRI slice generator is then trained and connected to the platform for inference. In the three-dimensional stage, 180 continuous MRI slices are used as supervision, and a 300000-point initial point cloud serves as the spatial starting point. After 30000 iterations of 3DGS training, the system obtains a representation with 556424 Gaussian points, and a 60000-point preview cloud is exported for frontend interaction. The engineering platform uses a FastAPI and React architecture. The backend manages data, generation, analysis, task, and 3D asset APIs, while the frontend provides the dashboard, generation workbench, medical image viewer, image analysis report, trust evaluation, point cloud display, and AI assistant. Experimental records show that the trained StyleGAN2 model can generate MRI-like slices with recognizable brain grayscale structures. At iteration 30000, 3DGS training reaches an L1 loss of 0.0459 and a PSNR of 20.0909 dB. These metrics mainly reflect rendering reconstruction consistency in this experiment. Project logs, model snapshots, rendered slices, metric curves, and frontend display assets together show that the platform provides a reproducible workflow from MRI data processing to two-dimensional generation and three-dimensional Gaussian representation.
\end{thesisenabstract}

\begin{thesisenkeywords}
\noindent{\bfseries Key Words}{: MRI Medical Image; 3D Gaussian Splatting; StyleGAN2; Medical Image Pseudo-Generation; Teaching Platform}
\end{thesisenkeywords}
